% Original topic-level problem, solution and MCQ reasoning.
\subsection{বাস্তব সমস্যা: কর্মী, Supplier ও Customer}
\chTwoPair{\textbf{সমস্যা:} একটি প্রতিষ্ঠানের কর্মীরা বেতনের বিবরণী দেখবেন। অনুমোদিত সরবরাহকারীরা অর্ডারের অবস্থা জানবেন এবং গ্রাহকেরা সাধারণ পণ্যের পৃষ্ঠা দেখবেন। সব তথ্য সবার জন্য উন্মুক্ত রাখা যাবে না।\par
\textbf{সমাধান:} কর্মীদের সীমিত অভ্যন্তরীণ সেবা ইন্ট্রানেটে এবং সরবরাহকারীদের নিয়ন্ত্রিত বহিরাগত প্রবেশ এক্সট্রানেটে রাখা যায়। সাধারণ পৃষ্ঠাটি ইন্টারনেটে প্রকাশ করা যায়। প্রতিটি ক্ষেত্রে উপযুক্ত অনুমতি ও পরিচয় যাচাই প্রয়োজন।}{\textbf{ব্যবহার:} ব্যবসায় অনুমোদিত তথ্য ভাগাভাগি, কলেজে কর্মীদের পোর্টাল, সরবরাহকারীদের সঙ্গে সমন্বয় ও শাখার নেটওয়ার্কে এই ব্যবস্থা ব্যবহার করা যায়। একই প্রতিষ্ঠানের LAN/WAN-এ প্রবেশাধিকারের এসব ধরন থাকতে পারে।\par
\textbf{Area ও Access-এর পার্থক্য:} ভৌগোলিক বিস্তার এবং প্রবেশাধিকার আলাদা বিষয়। প্রথমটি নেটওয়ার্কের এলাকা, দ্বিতীয়টি ব্যবহারকারীর অনুমতি বোঝায়। দূরের শাখাও প্রতিষ্ঠানের সীমিত অভ্যন্তরীণ সেবা ব্যবহার করতে পারে।\par
\textbf{Hardware Resource ও Data Sharing:} প্রিন্টার হার্ডওয়্যার রিসোর্স। ফাইল বা ডেটা ভাগাভাগি হলো Data Sharing। দুটিই নেটওয়ার্কিংয়ের উদ্দেশ্য; তবে উদাহরণ ও ভাগাভাগি করা সম্পদ আলাদা।\par}

\textbf{পরীক্ষায় প্রচলিত দূরত্বের উদাহরণ:} প্রচলিত তুলনায় PAN প্রায় ১০ m, LAN প্রায় ১ km, CAN প্রায় ১–৫ km এবং MAN প্রায় ৫০ km বিস্তৃত বলে উদাহরণ দেওয়া হয়। এগুলো শ্রেণি বোঝানোর আনুমানিক মান। কেবলের মান অনুযায়ী সর্বোচ্চ দৈর্ঘ্য বা প্রতিটি নেটওয়ার্কের অপরিবর্তনীয় ব্যাসার্ধ নয়। একই ক্যাম্পাসের একাধিক ভবনের নেটওয়ার্ক CAN হতে পারে; আবার পাশাপাশি কয়েকটি ভবন সাধারণ LAN-এর অংশও হতে পারে। প্রশ্নের পরিসর ও বিকল্প বিবেচনা করতে হয়।

\textbf{একই কলেজের নেটওয়ার্কের তিনটি ভিত্তি:} ক্যাম্পাসের একাধিক LAN যুক্ত থাকলে CAN-এর ধারণা প্রযোজ্য; প্রশ্নের পরিসর বিবেচনা করতে হয়। ফলাফলের সেবা সার্ভার থেকে দেওয়া হলে Client--server Architecture প্রযোজ্য। সীমিত অভ্যন্তরীণ পোর্টাল হলে সেটি Intranet। ভৌগোলিক বিস্তৃতি, সেবার আর্কিটেকচার ও প্রবেশাধিকার আলাদা ভিত্তি। একই নেটওয়ার্ককে এই তিন দিক থেকে বর্ণনা করা যায়।\par
